\documentclass[10pt,a4paper]{amsart}
\usepackage{amsmath,amssymb,amsthm,graphicx}
\usepackage{verbatim}
\newenvironment{reference}{\comment}{\endcomment}
%\newenvironment{reference}{}{}
\newenvironment{slogan}{\comment}{\endcomment}
\newenvironment{history}{\comment}{\endcomment}

% For commutative diagrams we use Xy-pic
\usepackage[all]{xy}

% We use 2cell for 2-commutative diagrams.
\xyoption{2cell}
% Theorem environments.
%
\theoremstyle{plain}
\newtheorem{theorem}[subsection]{Theorem}
\newtheorem{proposition}[subsection]{Proposition}
\newtheorem{lemma}[subsection]{Lemma}

\theoremstyle{definition}
\newtheorem{definition}[subsection]{Definition}
\newtheorem{example}[subsection]{Example}
\newtheorem{exercise}[subsection]{Exercise}
\newtheorem{situation}[subsection]{Situation}

\theoremstyle{remark}
\newtheorem{remark}[subsection]{Remark}
\newtheorem{remarks}[subsection]{Remarks}

\numberwithin{equation}{subsection}

% Macros
%
\def\lim{\mathop{\mathrm{lim}}\nolimits}
\def\colim{\mathop{\mathrm{colim}}\nolimits}
\def\Spec{\mathop{\mathrm{Spec}}}
\def\Hom{\mathop{\mathrm{Hom}}\nolimits}
\def\Ext{\mathop{\mathrm{Ext}}\nolimits}
\def\SheafHom{\mathop{\mathcal{H}\!\mathit{om}}\nolimits}
\def\SheafExt{\mathop{\mathcal{E}\!\mathit{xt}}\nolimits}
\def\Sch{\mathit{Sch}}
\def\Mor{\mathop{\mathrm{Mor}}\nolimits}
\def\Ob{\mathop{\mathrm{Ob}}\nolimits}
\def\Sh{\mathop{\mathit{Sh}}\nolimits}
\def\NL{\mathop{N\!L}\nolimits}
\def\CH{\mathop{\mathrm{CH}}\nolimits}
\def\proetale{{pro\text{-}\acute{e}tale}}
\def\etale{{\acute{e}tale}}
\def\QCoh{\mathit{QCoh}}
\def\Ker{\mathop{\mathrm{Ker}}}
\def\Im{\mathop{\mathrm{Im}}}
\def\Coker{\mathop{\mathrm{Coker}}}
\def\Coim{\mathop{\mathrm{Coim}}}

% Boxtimes
%
\DeclareMathSymbol{\boxtimes}{\mathbin}{AMSa}{"02}

%
% Macros for moduli stacks/spaces
%
\def\QCohstack{\mathcal{QC}\!\mathit{oh}}
\def\Cohstack{\mathcal{C}\!\mathit{oh}}
\def\Spacesstack{\mathcal{S}\!\mathit{paces}}
\def\Quotfunctor{\mathrm{Quot}}
\def\Hilbfunctor{\mathrm{Hilb}}
\def\Curvesstack{\mathcal{C}\!\mathit{urves}}
\def\Polarizedstack{\mathcal{P}\!\mathit{olarized}}
\def\Complexesstack{\mathcal{C}\!\mathit{omplexes}}
% \Pic is the operator that assigns to X its picard group, usage \Pic(X)
% \Picardstack_{X/B} denotes the Picard stack of X over B
% \Picardfunctor_{X/B} denotes the Picard functor of X over B
\def\Pic{\mathop{\mathrm{Pic}}\nolimits}
\def\Picardstack{\mathcal{P}\!\mathit{ic}}
\def\Picardfunctor{\mathrm{Pic}}
\def\Deformationcategory{\mathcal{D}\!\mathit{ef}}
\usepackage{cleveref}
\crefname{originallemma}{original lemma}{original lemmas}
\Crefname{originallemma}{Original lemma}{Original lemmas}
\makeatletter
\newlabel{categories-lemma-yoneda}{{001P}{source}{}{sourcepointer-1}{}}
\newlabel{categories-lemma-yoneda@cref}{{[originallemma][0][]001P}{[1][0][]source}{}{}{}}
\makeatother
\crefname{originallemma}{original lemma}{original lemmas}
\Crefname{originallemma}{Original lemma}{Original lemmas}
\makeatletter
\newlabel{cohomology-lemma-apply-Leray}{{01F4}{source}{}{sourcepointer-2}{}}
\newlabel{cohomology-lemma-apply-Leray@cref}{{[originallemma][0][]01F4}{[1][0][]source}{}{}{}}
\makeatother
\crefname{originallemma}{original lemma}{original lemmas}
\Crefname{originallemma}{Original lemma}{Original lemmas}
\makeatletter
\newlabel{cohomology-lemma-projection-formula}{{01E8}{source}{}{sourcepointer-3}{}}
\newlabel{cohomology-lemma-projection-formula@cref}{{[originallemma][0][]01E8}{[1][0][]source}{}{}{}}
\makeatother
\crefname{originallemma}{original lemma}{original lemmas}
\Crefname{originallemma}{Original lemma}{Original lemmas}
\makeatletter
\newlabel{duality-lemma-dualizing-module-proper-over-A}{{0AWP}{source}{}{sourcepointer-4}{}}
\newlabel{duality-lemma-dualizing-module-proper-over-A@cref}{{[originallemma][0][]0AWP}{[1][0][]source}{}{}{}}
\makeatother
\crefname{originalreference}{original reference}{original references}
\Crefname{originalreference}{Original reference}{Original references}
\makeatletter
\newlabel{example-rational-bridge}{{0E3M}{source}{}{sourcepointer-5}{}}
\newlabel{example-rational-bridge@cref}{{[originalreference][0][]0E3M}{[1][0][]source}{}{}{}}
\makeatother
\crefname{originalreference}{original reference}{original references}
\Crefname{originalreference}{Original reference}{Original references}
\makeatletter
\newlabel{lemma-closed-subscheme-reduced-gorenstein}{{0E34}{source}{}{sourcepointer-6}{}}
\newlabel{lemma-closed-subscheme-reduced-gorenstein@cref}{{[originalreference][0][]0E34}{[1][0][]source}{}{}{}}
\makeatother
\crefname{originalreference}{original reference}{original references}
\Crefname{originalreference}{Original reference}{Original references}
\makeatletter
\newlabel{lemma-duality-dim-1-CM}{{0BS3}{source}{}{sourcepointer-7}{}}
\newlabel{lemma-duality-dim-1-CM@cref}{{[originalreference][0][]0BS3}{[1][0][]source}{}{}{}}
\makeatother
\crefname{originalreference}{original reference}{original references}
\Crefname{originalreference}{Original reference}{Original references}
\makeatletter
\newlabel{lemma-genus-zero-pic}{{0C6M}{source}{}{sourcepointer-8}{}}
\newlabel{lemma-genus-zero-pic@cref}{{[originalreference][0][]0C6M}{[1][0][]source}{}{}{}}
\makeatother
\crefname{originallemma}{original lemma}{original lemmas}
\Crefname{originallemma}{Original lemma}{Original lemmas}
\makeatletter
\newlabel{more-morphisms-lemma-bijection-on-Pic}{{0E24}{source}{}{sourcepointer-9}{}}
\newlabel{more-morphisms-lemma-bijection-on-Pic@cref}{{[originallemma][0][]0E24}{[1][0][]source}{}{}{}}
\makeatother
\title{The contraction of a rational bridge pulls back the dualizing sheaf isomorphically}
\author{The Stacks Project Authors}
\date{stacks-project:curves.tex:local-snapshot-20261012; GNU FDL 1.2 or later}
\begin{document}
\maketitle
\noindent\textit{Editorial note.} All delivered mathematical text is the author's own, copied from the pinned Stacks Project chapter file curves.tex; no statement or proof step was rewritten, reordered or removed.\par\medskip
\noindent\textit{Editorial note.} The counted proof is the author's complete argument: the category equivalence with its quasi-inverse, the restriction sequence on the bridge with the author's small-detail remark on the degree count, the vanishing input, the commuting square of functors, and the Yoneda conclusion are all retained verbatim.\par\medskip
\noindent\textit{Editorial note.} The label lemma-rational-bridge-canonical names the counted statement, and the proof environment immediately following it in the source is delivered verbatim as its proof.\par\medskip
\noindent\textit{Editorial note.} The counted statement is carried by the author's own lemma environment in the chapter source.\par\medskip
\noindent\textit{Allowed standard prerequisites (editorial).}\begin{itemize}
\item Dualizing module by representability. Same chapter, Lemma lemma-duality-dim-1-CM; Duality for Schemes, Lemma duality-lemma-dualizing-module-proper-over-A. Characterizes omega\_X as the coherent module representing F -> Hom\_k(H\textasciicircum{}1(X, F), k), the representability framework in which the conclusion is drawn.
\item Picard behaviour under the contraction. More on Morphisms, Lemma more-morphisms-lemma-bijection-on-Pic; Cohomology, Lemmas cohomology-lemma-projection-formula and cohomology-lemma-apply-Leray. Gives essential surjectivity of c\textasciicircum{}* on the relevant subcategories, the projection formula c\_*c\textasciicircum{}*N = N, R\textasciicircum{}1c\_*c\textasciicircum{}*N = 0, and the Leray comparison of the two H\textasciicircum{}1-dual functors.
\end{itemize}\medskip
\ifdefined\hypertarget\hypertarget{sourcepointer-1}{}\fi
\noindent\textit{Source reference.} Original lemma 001P: sources/stacks-chow/tags.txt, source lines 70--70. Prerequisite pointer invoked at proof line(s) 6429; it names a standard fact of a sibling chapter that the counted argument uses, and every step of the counted proof built on it is retained verbatim in the delivered body.\par\medskip
\ifdefined\hypertarget\hypertarget{sourcepointer-2}{}\fi
\noindent\textit{Source reference.} Original lemma 01F4: sources/stacks-chow/tags.txt, source lines 1857--1857. Prerequisite pointer invoked at proof line(s) 6416; it names a standard fact of a sibling chapter that the counted argument uses, and every step of the counted proof built on it is retained verbatim in the delivered body.\par\medskip
\ifdefined\hypertarget\hypertarget{sourcepointer-3}{}\fi
\noindent\textit{Source reference.} Original lemma 01E8: sources/stacks-chow/tags.txt, source lines 1825--1825. Prerequisite pointer invoked at proof line(s) 6397; it names a standard fact of a sibling chapter that the counted argument uses, and every step of the counted proof built on it is retained verbatim in the delivered body.\par\medskip
\ifdefined\hypertarget\hypertarget{sourcepointer-4}{}\fi
\noindent\textit{Source reference.} Original lemma 0AWP: sources/stacks-chow/tags.txt, source lines 13727--13727. Prerequisite pointer invoked at proof line(s) 6382; it names a standard fact of a sibling chapter that the counted argument uses, and every step of the counted proof built on it is retained verbatim in the delivered body.\par\medskip
\ifdefined\hypertarget\hypertarget{sourcepointer-5}{}\fi
\noindent\textit{Source reference.} Original reference 0E3M: sources/stacks-curves/pinned-original-source.tex, source lines 6201--6201. Prerequisite pointer invoked at proof line(s) 6369; it names a standard fact of the same chapter that the counted argument uses, and every step of the counted proof built on it is retained verbatim in the delivered body.\par\medskip
\ifdefined\hypertarget\hypertarget{sourcepointer-6}{}\fi
\noindent\textit{Source reference.} Original reference 0E34: sources/stacks-curves/pinned-original-source.tex, source lines 1012--1012. Prerequisite pointer invoked at proof line(s) 6407; it names a standard fact of the same chapter that the counted argument uses, and every step of the counted proof built on it is retained verbatim in the delivered body.\par\medskip
\ifdefined\hypertarget\hypertarget{sourcepointer-7}{}\fi
\noindent\textit{Source reference.} Original reference 0BS3: sources/stacks-curves/pinned-original-source.tex, source lines 864--864. Prerequisite pointer invoked at proof line(s) 6380; it names a standard fact of the same chapter that the counted argument uses, and every step of the counted proof built on it is retained verbatim in the delivered body.\par\medskip
\ifdefined\hypertarget\hypertarget{sourcepointer-8}{}\fi
\noindent\textit{Source reference.} Original reference 0C6M: sources/stacks-curves/pinned-original-source.tex, source lines 2005--2005. Prerequisite pointer invoked at proof line(s) 6409; it names a standard fact of the same chapter that the counted argument uses, and every step of the counted proof built on it is retained verbatim in the delivered body.\par\medskip
\ifdefined\hypertarget\hypertarget{sourcepointer-9}{}\fi
\noindent\textit{Source reference.} Original lemma 0E24: sources/stacks-chow/tags.txt, source lines 17594--17594. Prerequisite pointer invoked at proof line(s) 6395; it names a standard fact of a sibling chapter that the counted argument uses, and every step of the counted proof built on it is retained verbatim in the delivered body.\par\medskip
\begin{lemma}
\label{lemma-rational-bridge-canonical}
Let $c : X \to Y$ be the contraction of a rational bridge
(Example \ref{example-rational-bridge}).
Then $c^*\omega_Y \cong \omega_X$.
\end{lemma}

\begin{proof}
You can prove this by direct computation, but we prefer to use the
characterization of $\omega_X$ as the coherent $\mathcal{O}_X$-module
which represents the functor
$\textit{Coh}(\mathcal{O}_X) \to \textit{Sets}$,
$\mathcal{F} \mapsto \Hom_k(H^1(X, \mathcal{F}), k) =
H^1(X, \mathcal{F})^\vee$, see
Lemma \ref{lemma-duality-dim-1-CM} or
Duality for Schemes, Lemma
\ref{duality-lemma-dualizing-module-proper-over-A}.

\medskip\noindent
To be precise, denote $\mathcal{C}_Y$ the category whose objects are
invertible $\mathcal{O}_Y$-modules and whose maps are
$\mathcal{O}_Y$-module homomorphisms. Denote $\mathcal{C}_X$ the category
whose objects are invertible $\mathcal{O}_X$-modules $\mathcal{L}$ with
$\mathcal{L}|_C \cong \mathcal{O}_C$ and whose maps are
$\mathcal{O}_Y$-module homomorphisms. We claim that the functor
$$
c^* : \mathcal{C}_Y \to \mathcal{C}_X
$$
is an equivalence of categories. Namely, by
More on Morphisms, Lemma \ref{more-morphisms-lemma-bijection-on-Pic}
it is essentially surjective. Then the projection formula
(Cohomology, Lemma \ref{cohomology-lemma-projection-formula})
shows $c_*c^*\mathcal{N} = \mathcal{N}$ and hence $c^*$
is an equivalence with quasi-inverse given by $c_*$.

\medskip\noindent
We claim $\omega_X$ is an object of $\mathcal{C}_X$. Namely, we have a
short exact sequence
$$
0 \to \omega_C \to \omega_X|_C \to \mathcal{O}_{C \cap X'} \to 0
$$
See Lemma \ref{lemma-closed-subscheme-reduced-gorenstein}.
Taking degrees we find $\deg(\omega_X|_C) = 0$ (small detail omitted).
Thus $\omega_X|_C$ is trivial by Lemma \ref{lemma-genus-zero-pic}
and $\omega_X$ is an object of $\mathcal{C}_X$.

\medskip\noindent
Since $R^1c_*\mathcal{O}_X = 0$ the projection formula shows that
$R^1c_*c^*\mathcal{N} = 0$ for $\mathcal{N} \in \Ob(\mathcal{C}_Y)$.
Therefore the Leray spectral sequence
(Cohomology, Lemma \ref{cohomology-lemma-apply-Leray})
the diagram
$$
\xymatrix{
\mathcal{C}_Y \ar[rr]_{c^*} \ar[dr]_{H^1(Y, -)^\vee} & &
\mathcal{C}_X \ar[ld]^{H^1(X, -)^\vee} \\
& \textit{Sets}
}
$$
of categories and functors is commutative. Since
$\omega_Y \in \Ob(\mathcal{C}_Y)$ represents the south-east arrow and
$\omega_X \in \Ob(\mathcal{C}_X)$ represents the south-east arrow
we conclude by the Yoneda lemma
(Categories, Lemma \ref{categories-lemma-yoneda}).
\end{proof}

\end{document}
